% ═══════════════════════════════════════════════════════════════
% SPANISCH LEHRERHINWEISE - Sequenz 07, Block b: Números 21-100
% ═══════════════════════════════════════════════════════════════
% Fachfarbe: #c41e3a (Karminrot)
% ═══════════════════════════════════════════════════════════════

\documentclass[11pt, a4paper]{article}

% ═══════════════════════════════════════════════════════════════
% SW-MODUS TOGGLE
% ═══════════════════════════════════════════════════════════════
\newif\ifswmodus
\swmodusfalse    % Standard: Farbe

% ═══════════════════════════════════════════════════════════════
% PAKETE
% ═══════════════════════════════════════════════════════════════

\usepackage[utf8]{inputenc}
\usepackage[T1]{fontenc}
\usepackage[ngerman]{babel}
\usepackage{helvet}
\renewcommand{\familydefault}{\sfdefault}

\usepackage[margin=2cm]{geometry}
\usepackage{enumitem}
\usepackage{tabularx}
\usepackage{booktabs}
\usepackage{longtable}

\usepackage{xcolor}
\usepackage{amssymb}
\usepackage{tcolorbox}
\tcbuselibrary{skins}

\usepackage{fancyhdr}
\usepackage{lastpage}
\usepackage{needspace}

% ═══════════════════════════════════════════════════════════════
% FARBEN
% ═══════════════════════════════════════════════════════════════

\definecolor{spanisch-color}{HTML}{c41e3a}
\definecolor{grau-color}{HTML}{888888}
\definecolor{hellgrau-color}{HTML}{f5f5f5}
\definecolor{basis-color}{HTML}{2a7a4b}
\definecolor{standard-color}{HTML}{2c5aa0}
\definecolor{erweiterung-color}{HTML}{7b1fa2}
\definecolor{loesung-color}{HTML}{c62828}

\definecolor{sw-dark}{HTML}{000000}
\definecolor{sw-gray}{HTML}{666666}
\definecolor{sw-white}{HTML}{ffffff}

\ifswmodus
    \colorlet{spanisch}{sw-dark}
    \colorlet{grau}{sw-gray}
    \colorlet{hellgrau}{sw-white}
    \colorlet{basis}{sw-dark}
    \colorlet{standard}{sw-dark}
    \colorlet{erweiterung}{sw-dark}
    \colorlet{loesung}{sw-dark}
\else
    \colorlet{spanisch}{spanisch-color}
    \colorlet{grau}{grau-color}
    \colorlet{hellgrau}{hellgrau-color}
    \colorlet{basis}{basis-color}
    \colorlet{standard}{standard-color}
    \colorlet{erweiterung}{erweiterung-color}
    \colorlet{loesung}{loesung-color}
\fi

\newcommand{\fachfarbe}{spanisch}

% ═══════════════════════════════════════════════════════════════
% VARIABLEN
% ═══════════════════════════════════════════════════════════════

\newcommand{\thema}{Números 21-100 y precios}
\newcommand{\sequenz}{07}
\newcommand{\block}{b}
\newcommand{\fach}{Spanisch}
\newcommand{\klasse}{7}
\newcommand{\dauer}{90 Min. (Doppelstunde)}
\newcommand{\lerngruppengroesse}{24}
\newcommand{\datum}{\today}

% ═══════════════════════════════════════════════════════════════
% KOPF- UND FUSSZEILE
% ═══════════════════════════════════════════════════════════════

\pagestyle{fancy}
\fancyhf{}
\renewcommand{\headrulewidth}{0.4pt}
\renewcommand{\footrulewidth}{0.4pt}

\lhead{\fach{} -- Klasse \klasse}
\chead{\textbf{LH-\sequenz\block{} (Lehrerhinweise)}}
\rhead{\datum}

\lfoot{}
\cfoot{}
\rfoot{Seite \thepage\ von \pageref{LastPage}}

% ═══════════════════════════════════════════════════════════════
% BEFEHLE
% ═══════════════════════════════════════════════════════════════

\newcommand{\B}{\textcolor{basis}{\textbf{[B]}}}
\newcommand{\Std}{\textcolor{standard}{\textbf{[S]}}}
\newcommand{\E}{\textcolor{erweiterung}{\textbf{[E]}}}

\newcommand{\loesung}[1]{\textcolor{loesung}{\textbf{#1}}}

\newcommand{\es}[1]{\textcolor{\fachfarbe}{\textbf{#1}}}

\newtcolorbox{kurzplan}{
    colback=hellgrau,
    colframe=\fachfarbe,
    colbacktitle=white,
    coltitle=\fachfarbe,
    boxrule=1pt,
    arc=0pt,
    fonttitle=\bfseries,
    attach boxed title to top left={yshift=-2mm, xshift=5mm},
    boxed title style={boxrule=0pt, colback=white},
    title=Kurzplan
}

\newtcolorbox{differenzierung}{
    colback=white,
    colframe=grau,
    colbacktitle=white,
    coltitle=grau,
    boxrule=0.5pt,
    arc=0pt,
    fonttitle=\bfseries,
    attach boxed title to top left={yshift=-2mm, xshift=5mm},
    boxed title style={boxrule=0pt, colback=white},
    title=Differenzierung
}

\newtcolorbox{fehlerbox}{
    colback=white,
    colframe=loesung,
    colbacktitle=white,
    coltitle=loesung,
    boxrule=0.5pt,
    arc=0pt,
    fonttitle=\bfseries,
    attach boxed title to top left={yshift=-2mm, xshift=5mm},
    boxed title style={boxrule=0pt, colback=white},
    title=Häufige Fehler
}

\newtcolorbox{materialbox}{
    colback=white,
    colframe=\fachfarbe,
    colbacktitle=white,
    coltitle=\fachfarbe,
    boxrule=0.5pt,
    arc=0pt,
    fonttitle=\bfseries,
    attach boxed title to top left={yshift=-2mm, xshift=5mm},
    boxed title style={boxrule=0pt, colback=white},
    title=Material-Checkliste
}

% ═══════════════════════════════════════════════════════════════
% DOKUMENT
% ═══════════════════════════════════════════════════════════════

\begin{document}

% ═══════════════════════════════════════════════════════════════
% KOPFBEREICH
% ═══════════════════════════════════════════════════════════════

\begin{center}
    {\Large\bfseries\textcolor{\fachfarbe}{Lehrerhinweise: \thema}}\\[0.3em]
    {\normalsize LH-\sequenz\block{} | \dauer}
\end{center}

\vspace{10pt}

% ═══════════════════════════════════════════════════════════════
% 1. KURZPLAN
% ═══════════════════════════════════════════════════════════════

\section*{1. Stundenverlauf (Kurzplan)}

\textbf{1. Stunde (45 Min.)}

\begin{tabularx}{\textwidth}{|c|l|X|l|}
\hline
\textbf{Zeit} & \textbf{Phase} & \textbf{Inhalt} & \textbf{Material} \\
\hline
0--5 & Einstieg & Blitzrunde: Zahlen 0-20 wiederholen & -- \\
\hline
5--15 & Input 1 & Zahlen 21-29 einführen (ein Wort!) & Tafel \\
\hline
15--22 & Übung 1 & Chorisches Sprechen: 21-29 & -- \\
\hline
22--30 & Input 2 & Zehner 30-100 einführen, Systematik & Tafel \\
\hline
30--40 & Übung 2 & Zahlen-Bingo & Bingo-Karten \\
\hline
40--45 & Sicherung & Schwierige Zahlen wiederholen & -- \\
\hline
\end{tabularx}

\vspace{1em}

\textbf{2. Stunde (45 Min.)}

\begin{tabularx}{\textwidth}{|c|l|X|l|}
\hline
\textbf{Zeit} & \textbf{Phase} & \textbf{Inhalt} & \textbf{Material} \\
\hline
0--5 & Wiederholung & Quiz: ``¿Qué número es?'' & Karten \\
\hline
5--12 & Input 3 & Preise: ``Cuesta X euros (y Y céntimos)'' & Preisschilder \\
\hline
12--25 & Anwendung & Einkaufsdialoge (PA) & Preisschilder \\
\hline
25--35 & Festigung & AB-07b bearbeiten & AB \\
\hline
35--40 & Sicherung & Lösungen vergleichen & Tafel \\
\hline
40--45 & Abschluss & Zusammenfassung, HA, LP-Hinweis & -- \\
\hline
\end{tabularx}

% ═══════════════════════════════════════════════════════════════
% 2. DIFFERENZIERUNG
% ═══════════════════════════════════════════════════════════════

\section*{2. Differenzierung}

\begin{differenzierung}
\textbf{Legende:} \B{} Basis (BR) | \Std{} Standard (MR) | \E{} Erweiterung (GY)

\vspace{0.5em}

\textbf{WICHTIG:} Diese Marker sind NUR für die Lehrkraft. Auf Schülermaterial erscheinen KEINE Niveaumarkierungen!
\end{differenzierung}

\vspace{1em}

\begin{tabularx}{\textwidth}{|l|X|X|X|}
\hline
& \B{} \textbf{Basis} & \Std{} \textbf{Standard} & \E{} \textbf{Erweiterung} \\
\hline
\textbf{Aufg. 1} & Mit Zahlentabelle & Selbstständig & + Zusätzliche Zuordnung \\
\hline
\textbf{Aufg. 2} & Nur 2 Zahlen & Alle 4 Zahlen & + 2 eigene Zahlen \\
\hline
\textbf{Aufg. 3} & Nur a) + b) & Alle 4 Aufgaben & + Rückübersetzung \\
\hline
\textbf{Aufg. 4} & Mit Zahlentabelle & Alle 4 Aufgaben & + 2 eigene Preise \\
\hline
\textbf{Aufg. 5} & Mit Satzanfängen & Selbstständig & + Erweiterter Dialog \\
\hline
\end{tabularx}

\vspace{1em}

\textbf{Konkrete Hinweise:}
\begin{itemize}
    \item \B{} SuS mit LRS: Zeitzugabe (+10 Min.), Zahlentabelle als Hilfe
    \item \B{} DaZ-SuS: Zahlen 21-29 vorher üben, Visualisierung nutzen
    \item \E{} Leistungsstarke: Rechenaufgaben auf Spanisch (z.B. ``veinticinco más treinta'')
\end{itemize}

% ═══════════════════════════════════════════════════════════════
% 3. HÄUFIGE FEHLER
% ═══════════════════════════════════════════════════════════════

\section*{3. Häufige Fehler}

\begin{fehlerbox}
\begin{tabularx}{\textwidth}{|l|X|X|}
\hline
\textbf{Fehler} & \textbf{Beispiel} & \textbf{Reaktion} \\
\hline
21-29 getrennt & *``veinti uno'' & Zusammenschreibung betonen, Tafel \\
\hline
Akzente vergessen & *``veintidos'' statt \es{veintidós} & Übertrieben betonen: ``veintidÓS'' \\
\hline
40/50 verwechselt & *``cuarenta'' für 50 & Eselsbrücke: \textit{cin}cuenta = \textit{cin}co \\
\hline
``y'' vergessen & *``treinta uno'' & Rhythmisch üben: ``TRÉIN-ta Y Ú-no'' \\
\hline
``céntimos'' falsch & *``centimos'' & th-Laut üben, nicht wie ``Cent'' \\
\hline
100 = ``ciento'' & *``ciento euros'' & Nur alleinstehend \es{cien}, sonst ciento \\
\hline
\end{tabularx}
\end{fehlerbox}

% ═══════════════════════════════════════════════════════════════
% 4. MATERIAL-CHECKLISTE
% ═══════════════════════════════════════════════════════════════

\section*{4. Material-Checkliste}

\begin{materialbox}
\begin{itemize}[label=$\square$]
    \item AB-\sequenz\block.pdf (\lerngruppengroesse{} Kopien)
    \item ML-\sequenz\block.pdf (1× für Lehrkraft)
    \item Lehrwerk: Qué pasa 1, S. 48-49
    \item Bingo-Karten (25 verschiedene, Zahlen 21-100)
    \item Preisschilder (Bildkarten mit Produkten + Preisen)
    \item Tafelmarker / Kreide
    \item Optional: LP-\sequenz\block.html (Tablets/PCs)
\end{itemize}
\end{materialbox}

% ═══════════════════════════════════════════════════════════════
% 5. LÖSUNGEN
% ═══════════════════════════════════════════════════════════════

\section*{5. Lösungen AB-\sequenz\block}

\textbf{Merkkasten 1: Zahlen 21-29}

\begin{tabular}{lll}
21 = \loesung{veintiuno} & 22 = \loesung{veintidós} & 23 = \loesung{veintitrés} \\
24 = \loesung{veinticuatro} & 25 = \loesung{veinticinco} & 26 = \loesung{veintiséis} \\
27 = \loesung{veintisiete} & 28 = \loesung{veintiocho} & 29 = \loesung{veintinueve} \\
\end{tabular}

\vspace{1em}

\textbf{Merkkasten 2: Zehner 30-100}

\begin{tabular}{llll}
30 = \loesung{treinta} & 40 = \loesung{cuarenta} & 50 = \loesung{cincuenta} & 60 = \loesung{sesenta} \\
70 = \loesung{setenta} & 80 = \loesung{ochenta} & 90 = \loesung{noventa} & 100 = \loesung{cien} \\
\end{tabular}

\vspace{1em}

\textbf{Merkkasten 3: Preise}

\begin{tabular}{ll}
¿Cuánto cuesta? & $\rightarrow$ \loesung{Wieviel kostet es?} \\
Cuesta X euros. & $\rightarrow$ \loesung{Es kostet X Euro.} \\
Cuesta X euros y Y céntimos. & $\rightarrow$ \loesung{Es kostet X Euro und Y Cent.} \\
\end{tabular}

\vspace{1em}

\textbf{Aufgabe 1: Zuordnung} (4 Punkte)

\begin{tabular}{ll}
25 & $\rightarrow$ \loesung{veinticinco} \\
47 & $\rightarrow$ \loesung{cuarenta y siete} \\
63 & $\rightarrow$ \loesung{sesenta y tres} \\
81 & $\rightarrow$ \loesung{ochenta y uno} \\
\end{tabular}

\vspace{1em}

\textbf{Aufgabe 2: Zahlen schreiben} (4 Punkte)

\begin{tabular}{ll}
a) 22 & $\rightarrow$ \loesung{veintidós} \\
b) 38 & $\rightarrow$ \loesung{treinta y ocho} \\
c) 54 & $\rightarrow$ \loesung{cincuenta y cuatro} \\
d) 99 & $\rightarrow$ \loesung{noventa y nueve} \\
\end{tabular}

\vspace{1em}

\textbf{Aufgabe 3: Preise lesen} (4 Punkte)

\begin{tabular}{ll}
a) & \loesung{23 €} \\
b) & \loesung{45 €} \\
c) & \loesung{78,50 €} \\
d) & \loesung{99 €} \\
\end{tabular}

\vspace{1em}

\textbf{Aufgabe 4: Preise sagen} (4 Punkte)

\begin{tabular}{ll}
a) 35 € & $\rightarrow$ \loesung{treinta y cinco} \\
b) 67 € & $\rightarrow$ \loesung{sesenta y siete} \\
c) 82,50 € & $\rightarrow$ \loesung{ochenta y dos ... cincuenta} \\
d) 100 € & $\rightarrow$ \loesung{cien} \\
\end{tabular}

\vspace{1em}

\textbf{Aufgabe 5: Dialog} (4 Punkte)

\begin{itemize}
    \item Lücke 1: \loesung{¿Cuánto cuesta} (1P)
    \item Lücke 2: \loesung{Cuesta [Preis] euros.} (1P)
    \item Lücke 3: \loesung{Aquí tiene.} (1P)
    \item Lücke 4: \loesung{Gracias} (1P)
\end{itemize}

% ═══════════════════════════════════════════════════════════════
% 6. HAUSAUFGABE
% ═══════════════════════════════════════════════════════════════

\section*{6. Hausaufgabe}

\begin{itemize}
    \item Lehrwerk S. 49, Aufgabe 3 (Zahlen schreiben)
    \item Vokabeln lernen: Zahlen 21-100, Preisausdrücke
    \item Optional: Lernpfad LP-\sequenz\block.html (digital)
\end{itemize}

% ═══════════════════════════════════════════════════════════════
% 7. REFLEXIONSNOTIZEN
% ═══════════════════════════════════════════════════════════════

\section*{7. Reflexionsnotizen (nach Durchführung)}

\begin{tabularx}{\textwidth}{|l|X|}
\hline
\textbf{Was lief gut?} & \\[2em]
\hline
\textbf{Was lief nicht?} & \\[2em]
\hline
\textbf{Zeitmanagement} & \\[2em]
\hline
\textbf{Für nächstes Mal} & \\[2em]
\hline
\end{tabularx}

\end{document}
